\begin{proposition}[G-15: Uniform third-jet dominance]\label{G-15}
For every $m\geq1$, the map
\[
 j_{\leq3}:\mathcal D_{m,4}\longrightarrow
   \bigoplus_{j=0}^{\min\{3,m\}}
          \operatorname{Sym}^{m-j}\mathbb C^3,
 \qquad \sum_{j=0}^m t^j q_j\longmapsto(q_0,q_1,q_2,q_3),
\]
with nonexistent components omitted, has Zariski-dense image.
For $m\geq4$, order three is the largest order for which the
full jet map can be dominant.

Consequently, if $n>m\geq1$, $a=n-m$, $d\geq1$, and
$\lambda\vdash nd$ has at most four rows, then
\[
 \lambda_1\leq ad+3
 \quad\Longrightarrow\quad
 m_\lambda(\ell^a f)\leq m_\lambda(\det_n)
\]
for every degree-$m$ form $f$ in the common ambient setting.
Every positive gap in this sector therefore satisfies
$\lambda_1\geq ad+4$. In particular a rectangular gap
with $\lambda=(nd/4)^4$ requires
\[
 d(4m-3n)>12.
\]
Since $4\mid nd$ for this rectangular shape, the last
quantity is divisible by four; equivalently,
$d(4m-3n)\geq16$ is necessary.
\end{proposition}

\paragraph{Background and proof strategy.}
We use the standard $\mathfrak{sl}_2$ decompositions
\[
 \operatorname{End}(\operatorname{Sym}^{m-1}\mathbb C^2)
       =\bigoplus_{l=0}^{m-1}\operatorname{Sym}^{2l}\mathbb C^2,
 \qquad
 \operatorname{Sym}^{r}(\operatorname{Sym}^2\mathbb C^2)
       =\bigoplus_{b=0}^{\lfloor r/2\rfloor}
             \operatorname{Sym}^{2r-4b}\mathbb C^2.
\]
These are the Clebsch--Gordan and harmonic decompositions;
see Fulton and Harris (1991), section~11.2, Exercises~11.11
and 11.14, printed pp.~151 and 153. We specify their
covariants, the finite orthogonal-polynomial calculation,
and the required rank argument below. No generic strong
Lefschetz property or Gulliksen--Neg\aa rd theorem is
assumed in this proof.

\begin{proof}
The case $m=1$ is immediate. Fix $m\geq2$ and let
$\rho$ be the irreducible representation of
$\mathfrak{sl}_2$ on $V=\operatorname{Sym}^{m-1}\mathbb C^2$.
Use standard generators $H,E,F$, with
$[H,E]=2E$, $[H,F]=-2F$, and $[E,F]=H$.
Choose the invariant Hermitian metric for which
$\rho(H)$ is Hermitian and $\rho(E)^*=\rho(F)$.
Its $m$ distinct $H$-eigenvalues are
\[
 \lambda_j=m-1-2j,\qquad 0\leq j<m.
\]
For $z=(z_1,z_2,z_3)$ put
\[
 A(z)=z_1H+z_2(E+F)+\mathrm i z_3(E-F),\qquad
 q=z_1^2+z_2^2+z_3^2,
\]
and consider the explicit linear pencil
\begin{equation}\label{G15:pencil}
 M(t,z)=tI_m-\rho(A(z)).
\end{equation}
The invariant quadratic form is normalized so that
$q(H)=1$. In the defining two-dimensional representation,
$A(z)^2=q(z)I_2$.

Let $\mathcal H_l$ be the space of harmonic homogeneous
polynomials of degree $l$ in $z_1,z_2,z_3$. It is the
irreducible representation of highest weight $2l$ and
dimension $2l+1$. The harmonic decomposition reads
\begin{equation}\label{G15:harmonic}
 \mathbb C[z_1,z_2,z_3]
     =\bigoplus_{l\geq0}\mathcal H_l\otimes\mathbb C[q].
\end{equation}
One can obtain this directly by applying the Laplacian to
$q^b h_l$: for harmonic $h_l$,
\[
 \Delta(q^b h_l)=2b(2l+2b+1)q^{b-1}h_l.
\]
This gives the direct decomposition and its dimensions in
characteristic zero. The indicated irreducibility is the
standard highest-weight calculation for the adjoint
three-dimensional representation.

\paragraph{1. Canonical harmonic matrix covariants.}
Write the multiplicity-free decomposition of
$\operatorname{End}(V)$ as
$\bigoplus_{l=0}^{m-1}\mathcal E_l$, where
$\mathcal E_l\simeq\mathcal H_l$, and let $\pi_l$ be its
equivariant projection. The trace Hermitian form makes
these distinct summands orthogonal. The polynomial map
\[
 A\longmapsto\pi_l\bigl(\rho(A)^l\bigr)
\]
is nonzero and has harmonic coefficient functions of
degree $l$. Indeed the symmetrized multiplication map
from the $l$th symmetric power of the adjoint
representation to $\operatorname{End}(V)$ is equivariant.
Its highest-weight vector $E^l$ maps to
$\rho(E)^l\ne0$ for $l<m$, and this vector has highest
weight $2l$, so lies in $\mathcal E_l$. The harmonic
decomposition says that this is the unique copy of
highest weight $2l$ in symmetric degree $l$.

The value of this covariant at $H$ is nonzero. Otherwise
equivariance and homogeneity would make it vanish on
every scalar multiple of the semisimple orbit of $H$,
which is the dense open set $q\ne0$, contradicting its
nonvanishing. This value is diagonal, since it has
Cartan weight zero.

These diagonal vectors have an elementary description.
Put the uniform positive inner product on functions on
$\{\lambda_0,\ldots,\lambda_{m-1}\}$:
\[
 \langle f,g\rangle=\frac1m\sum_{j=0}^{m-1}
                          f(\lambda_j)\overline{g(\lambda_j)}.
\]
Let $p_0,\ldots,p_{m-1}$ be the real orthonormal
polynomials obtained from $1,x,\ldots,x^{m-1}$, with
positive leading coefficients; thus $p_0=1$.
The zero-weight vector in $\mathcal E_l$ is proportional
to $p_l(\rho(H))$. To check the assertion, note that
$\rho(H)^r$ lies in the sum of the zero-weight spaces of
$\mathcal E_0,\ldots,\mathcal E_r$ by the symmetric-power
decomposition. The top projection is nonzero by the
preceding argument. Induction therefore identifies their
successive spans with those of $1,x,\ldots,x^r$ on the
spectrum. Orthogonality of the endomorphism summands
is precisely orthogonality for the displayed uniform
inner product. This proves the assertion, including
its degree and nonzero leading coefficient.

Normalize the covariant $T_l(A)$ so that
$T_l(H)=p_l(\rho(H))$. Its coefficient functions span
all of $\mathcal H_l$, since it is a nonzero map between
the corresponding irreducible representations.
Equivariance and \eqref{G15:harmonic} now give uniquely
\begin{equation}\label{G15:adjexpansion}
 \pi_l\operatorname{adj} M(t,z)=T_l(A(z))F_l(t,q),
 \qquad 0\leq l<m,
\end{equation}
where $F_l$ is a polynomial homogeneous of weighted
degree $m-1-l$ for $\deg t=1$, $\deg q=2$.

\paragraph{2. The scalar factors are trailing Jacobi determinants.}
Multiplication by $x$ on the functions on the $m$ nodes,
expressed in the orthonormal basis $p_0,\ldots,p_{m-1}$,
is a real symmetric tridiagonal matrix $J$.
Tridiagonality follows because $xp_l$ is orthogonal
to every $p_r$ with $r<l-1$. Its off-diagonal entries
$b_1,\ldots,b_{m-1}$ are strictly positive: they are
the ratios of the positive leading coefficients of
successive orthonormal polynomials. The nodes and
measure are symmetric under $x\mapsto-x$, so $p_l$
has parity $l$. Thus every diagonal entry
$\langle xp_l,p_l\rangle$ is zero.

Let $J^{[l+1]}$ denote the trailing principal submatrix
on indices $l+1,\ldots,m-1$, with the empty determinant
equal to one, and define
\[
 Q_l(t)=\det(tI-J^{[l+1]}),\qquad 0\leq l<m.
\]
The spectrum of $J$ is the node set. If
$D(t)=\prod_j(t-\lambda_j)$, evaluating
\eqref{G15:adjexpansion} at $H$ and taking its
orthonormal trace coefficient gives
\[
 F_l(t,1)=\frac1m\sum_jp_l(\lambda_j)
                         \frac{D(t)}{t-\lambda_j}.
\]
On the other hand the spectral expression for the
$(l,0)$ entry of $(tI-J)^{-1}$ is
\[
 \frac1m\sum_j\frac{p_l(\lambda_j)}{t-\lambda_j}.
\]
Expansion of the corresponding tridiagonal cofactor
shows that its numerator is
$b_1\cdots b_l Q_l(t)$: the unique path from index
zero to index $l$ contributes the indicated product,
and all remaining rows and columns give the trailing
principal determinant. Therefore
\begin{equation}\label{G15:jacobi}
 F_l(t,1)=b_1\cdots b_l Q_l(t).
\end{equation}
In particular every factor is nonzero. Weighted
homogeneity then determines $F_l(t,q)$ from this
univariate identity. All $Q_l$ have the parity of
their degree $m-1-l$.

We record the elementary spectral facts needed below.
Every nonempty trailing block is real symmetric,
tridiagonal, and has positive off-diagonal entries.
An eigenvector is determined by its first coordinate
using the three-term recurrence; that coordinate cannot
be zero. Hence all eigenvalues are simple. Conjugation
by the alternating-sign diagonal matrix negates the
block, so its eigenvalues occur in opposite pairs.
An even-sized block has no zero eigenvalue, as its
determinant at zero is a nonzero signed product of
alternate squared off-diagonal entries. An odd-sized
block has a simple zero eigenvalue. These assertions
also hold for blocks of sizes zero or one with the
obvious interpretations.

\paragraph{3. Cofactor multiples span every third jet.}
Let $S=\mathbb C[t,z_1,z_2,z_3]$, and let $\mathcal C_m$
be the degree-$m$ space spanned by all linear-form
multiples of the entries of $\operatorname{adj} M$.
It is the differential image of the determinant-pencil
map at \eqref{G15:pencil}. Formula
\eqref{G15:adjexpansion} implies that it contains
\[
 tF_l(t,q)\mathcal H_l\qquad(0\leq l<m).
\]
It also contains
\begin{equation}\label{G15:lowerproduct}
 qF_{l+1}(t,q)\mathcal H_l\qquad(0\leq l\leq m-2).
\end{equation}
For clarity, multiplication of a harmonic polynomial of
degree $l+1$ by a linear form has a nonzero projection
onto $q\mathcal H_l$. Indeed
$\Delta(z_i h)=2\partial_i h$; these derivatives span
$\mathcal H_l$ by equivariance and nonvanishing.
The space $\mathcal C_m$ is invariant, so projecting
its products onto this irreducible summand stays in
$\mathcal C_m$. Multiplication by the invariant
$F_{l+1}$ gives \eqref{G15:lowerproduct}.
Similarly the top harmonic component of
$\mathcal H_1\mathcal H_{m-1}$ supplies
$\mathcal H_m\subseteq\mathcal C_m$, since
$F_{m-1}$ is a nonzero constant.

The quotient $S_m/t^4S_{m-4}$ is exactly the third-jet
space, with $S_{m-4}=0$ if $m<4$. Its harmonic
decomposition has, for each $0\leq l\leq m$, the
multiplicity space
\begin{equation}\label{G15:targetspace}
 \operatorname{span}\{t^j q^{(D-j)/2}:
         0\leq j\leq\min\{3,D\},\ j\equiv D\pmod2\},
 \qquad D=m-l,
\end{equation}
tensored with $\mathcal H_l$. It has dimension one
for $D=0,1$, and dimension two for $D\geq2$.
The cases $D=0$ and $D=1$ are already supplied by
$\mathcal H_m$ and $tF_{m-1}\mathcal H_{m-1}$.

For $D\geq2$, specialize $q=1$. On the fixed
multiplicity space \eqref{G15:targetspace}, this is
an isomorphism onto the corresponding powers of $t$
modulo $t^4$. By \eqref{G15:jacobi}, it suffices to
show that $tQ_l(t)$ and $Q_{l+1}(t)$ are independent
modulo $t^4$.

If $D$ is even, $Q_l$ has odd degree $D-1$ and a
simple root at zero, while $Q_{l+1}$ has even degree
$D-2$ and nonzero value at zero. Thus $tQ_l$ starts
with a nonzero $t^2$ term, while $Q_{l+1}$ has a
nonzero constant term. They are independent in the
even two-dimensional space modulo $t^4$.

If $D$ is odd, then $D\geq3$. The block
$K=J^{[l+1]}$ is real symmetric and invertible,
of even size $D-1$. Write $e$ for its first
coordinate vector. The diagonal resolvent identity gives
\[
 \frac{Q_{l+1}(t)}{Q_l(t)}
    =e^{\mathsf T}(tI-K)^{-1}e.
\]
The alternating-sign symmetry fixes $e$ and negates
$K$, so all diagonal entries of odd powers of
$K^{-1}$ vanish. Expanding the inverse at zero gives
\[
 \frac{Q_{l+1}(t)}{tQ_l(t)}
   =-e^{\mathsf T}K^{-2}e
     -t^2 e^{\mathsf T}K^{-4}e
     +O(t^4).
\]
The coefficient of $t^2$ is strictly negative, since
$e^{\mathsf T}K^{-4}e=\|K^{-2}e\|^2>0$. Thus
$Q_{l+1}$ cannot be a constant multiple of $tQ_l$
modulo $t^4$. This proves independence in the odd
two-dimensional space, completing all cases of
\eqref{G15:targetspace}.

Consequently the differential image $\mathcal C_m$
surjects onto $S_m/t^4S_{m-4}$. The determinant-pencil
parameter map followed by third-jet projection has a
surjective differential at the explicit point
\eqref{G15:pencil}. In characteristic zero its image
is dense. This proves third-jet dominance for every
$m\geq2$, and hence for all $m\geq1$.

For $m\geq4$ the third-jet target dimension is
\[
 \sum_{j=0}^3\binom{m-j+2}{2}=2m^2+2.
\]
The fourth-jet target adds
$\binom{m-2}{2}>0$, whereas G-06 bounds
$\dim\mathcal D_{m,4}\leq2m^2+2$. The fourth-jet
map, and therefore every higher full jet map, cannot
be dominant. This establishes the maximal-order
assertion.

Finally apply the coefficient-weight argument of
G-14 with the now proved third-jet dominance. A
degree-$d$ coefficient monomial surviving on
$x_0^a q$ has $x_0$ weight
$ad+j_1+\cdots+j_d$. If $\lambda_1\leq ad+3$,
its evaluation depends only on the first four
coefficient forms $q_0,q_1,q_2,q_3$. Their dense
realization by determinants and block padding force
the same isotypic ideal containment as in G-14.
This proves the multiplicity inequality, including
the elementary cases $m<3$ where all available
jets are included. The rectangular inequality is
the equivalent condition
$nd/4>ad+3$, namely $d(4m-3n)>12$.
\end{proof}

\paragraph{What this resolves.}
The argument proves the needed rank for one explicit
pencil in every size and hence proves dominance. It
does not assume or assert a general strong Lefschetz
theorem for arbitrary Artinian determinantal algebras.
The positivity above is the squared norm for an
ordinary finite-dimensional real symmetric Jacobi
matrix; no arithmetic purity or RH input is used.
No literature-priority claim is made for this deduction.
